\begin{abstract}
本稿では，モデルの再学習を行わずに，学習時に無い物体を検出し，粗い語から細かい語までの確信度に応じた粒度で命名する粒度可変アノテーションを提案する．
提案手法は，COCOで学習したYOLO11xのボックスにOWLv2の物体らしさの上位の領域を加え，領域の特徴から求めた7806語の確率をWordNetの上位語へ集約し，集約した確率が閾値以上の間だけ上位語から下位語へ進む．
さらに，各アノテーションと経路上の各単語に，検証用の画像のみで14個の重みを求めた正しさの確率をスコアとして付け，利用者がスコアの閾値によって出力の数と粒度を選べるようにする．
LVISの人手の注釈を学習に用いていないOWLv2を用い，LVISのアノテーションでCOCO外の物体を評価した結果，COCO外の物体の検出は56.0\%，正しい上位語・下位語による命名は21.1\%であった．
COCO外の物体と対応した同じ領域に対し，正しく命名できた割合は，確率が最大の1語を出力する場合の25.9\%に対し提案手法で43.0\%であった．
提案したスコアは，検出のスコアのみを用いる場合と比べ，全ての誤りを数えるAPを0.171から0.197に高め，期待校正誤差を0.145から0.028に下げた．
一方，LVISのカテゴリまで正しい出力のみを数えるAPでは1語を出力する方法が上回り，評価の語彙を事前に与えた場合には集約の効果は無かった．
また，加えた領域の出力の多くは物体の一部や背景に対するものであり，精度を制限する主な要因は命名よりも領域の追加にあることを示した．
予備実験として，COCOの80クラスの範囲では，切り抜き画像に対するCLIPとSigLIPのラベルの一致率は0.56〜0.58であり，COCOで学習した物体認識モデル自身の0.96〜0.97に大きく劣ることも示した．
\end{abstract}

\noindent
{\bf キーワード：}粒度可変アノテーション，WordNet，文章・画像埋め込みモデル，CLIP，SigLIP，YOLO，DETR，OWLv2，オープンボキャブラリ，校正
